\makeatletter
\def\input@path{{./}{../}{../../}{../../../}{../../../../}{preamble/}{../preamble/}{../../preamble/}{../../../preamble/}}
\makeatother
\input{preamble_exams.tex}
\graphicspath{{./}{../}{../../}{../../../}{../../../../}{preamble/}{../preamble/}{../../preamble/}{../../../preamble/}}
\setYColumnWidth{0.22}
\setlength{\columnsep}{7mm}
\setlength{\fboxsep}{2pt}
% Account for the shared heading's colour-box padding.
\let\ExamSubsectionBox\SubsectionBox
\RenewDocumentCommand{\SubsectionBox}{s m}{%
  \begingroup
  \setlength{\linewidth}{\dimexpr\linewidth-2\fboxsep\relax}%
  \IfBooleanTF{#1}{\ExamSubsectionBox*{#2}}{\ExamSubsectionBox{#2}}%
  \endgroup
}
\begin{document}
\noindent
\begin{tabularx}{\textwidth}{|C{2.8cm}|C{\dimexpr\textwidth-6cm-4\tabcolsep-4\arrayrulewidth-1pt\relax}|C{2.8cm}|}
\hline
\includegraphics[width=2.5cm]{logo.png} &
\CellCenter{\textbf{GLOBAL ECONOMICS}\par
\textbf{Grade 10}\par
\textbf{Catch ups exam}\par
\textbf{Decision-making under uncertainty}\par
\textbf{Teacher: Nicolás López Cuéllar}} &
\CellCenter{\textbf{1st TERM}\par\textbf{2026--2027}}\\\hline
\end{tabularx}

\noindent
\begin{tabular}{|Y|Z|}\hline
{\small\textbf{Learning objective:} Characterise economic decisions under uncertainty, organise payoff tables, calculate profits from revenues and costs, and justify choices using Maximax, Maximin, Minimax Regret, EMV/Bayes, and probability sensitivity analysis.}
&
{\footnotesize\textbf{Assessment criteria}\par
	\textbf{C1.} Characterises decision-making problems under uncertainty by identifying decision alternatives, states of nature, consequences, and the economic or financial context of the decision.\par
	\textbf{C2.} Organises decision problems using payoff tables, placing alternatives, states of nature, and payoffs in a clear structure for comparison.\par
	\textbf{C3.} Calculates profits by combining revenues with fixed and variable costs across different states of nature, determining the corresponding payoffs for each alternative, and organising the results in payoff tables for comparison.\par
	\textbf{C4.} Applies the Maximin and Maximax criteria to make optimistic or conservative decisions without probabilities, selecting the alternative with the highest possible payoff or the best worst-case, interpreting risks or limitations.\par
	\textbf{C5.} Constructs opportunity loss tables and applies the Minimax Regret criterion to select the alternative that minimises the maximum opportunity loss.\par
	\textbf{C6.} Calculates Expected Monetary Value using stated or estimated probabilities and applies the Bayes decision criterion to select the alternative with the highest expected payoff.\par
	\textbf{C7.} Uses subjective probabilities and sensitivity analysis to examine how changes in probability assumptions affect the preferred decision.}\\\hline
\end{tabular}

\StudentInfo \vspace{-3mm}

\begin{multicols}{2}
\raggedcolumns
\SubsectionBox*{Criteria assessment}
This exam assesses C1--C7. Each question and task shows its points. \textbf{Passing requires a minimum of 5 cumulative points per criterion, including points earned in previous learning activities}. Show all procedures; answers without the required procedure receive no credit. Partial credit may be awarded for a correct method with arithmetic errors. Missing names or late submissions incur a 1-point deduction per criterion.


\begin{center}
	\small\setlength{\tabcolsep}{5pt}
	\begin{tabular}{|c|c|c|c|c||c|c|c|c|}
		\hline
		\textbf{\textnumero}&\textbf{C1}&\textbf{C2}&\textbf{C3}&\textbf{C4}&
		\textbf{\textnumero}&\textbf{C5}&\textbf{C6}&\textbf{C7}\\\hline
		1&3&6&--&--&6&2&--&--\\\hline
		2&3&--&6&--&7&2&--&--\\\hline
		3&--&--&--&2&8&2&--&--\\\hline
		4&--&--&--&2&9&--&2&--\\\hline
		5&--&--&--&2&10&--&2&--\\\hline
		&&&&&11&--&2&--\\\hline
		&&&&&12&--&--&3\\\hline
		&&&&&13&--&--&3\\\hline
		\textbf{Tot.}&\textbf{6}&\textbf{6}&\textbf{6}&\textbf{6}&
		\textbf{Tot.}&\textbf{6}&\textbf{6}&\textbf{6}\\\hline
	\end{tabular}
\end{center}


\begin{minipage}{\linewidth}
\SubsectionBox{River Restoration Events \score{C1-\sml{3} C2-\sml{6}}}
An environmental association must choose a river cleanup or a wetland planting event before volunteer turnout is known. Turnout will be strong or modest. The river cleanup provides a net contribution of 12 USD per participant and a fixed donor contribution of 105 USD per event. Wetland planting provides 19 USD per participant and a fixed donor contribution of 80 USD per event. Strong turnout adds a further net contribution of 6 USD per participant to either event; modest turnout adds 2 USD per participant. The river cleanup attracts 32 participants with strong turnout or 19 with modest turnout; wetland planting attracts 26 or 16, respectively. All contributions are additive, and no other adjustments are needed.

Characterize the key elements of the decision-making situation \sscore{C1-3} 
and construct the payoff table. \sscore{C2-6}

\end{minipage}\par

\begin{minipage}{\linewidth}
\SubsectionBox{Guided Heritage Tours \score{C1-\sml{3} C3-\sml{6}}}
A tour company must choose walking tours or minibus tours before visitor bookings are known. Bookings will be plentiful or limited. Walking tours generate 92 ticket sales with plentiful bookings or 61 with limited bookings; minibus tours generate 73 or 42, respectively. Base revenue per ticket is 32 USD for walking tours and 57 USD for minibus tours. Plentiful bookings add a surcharge of 8 USD per ticket to either tour type; limited bookings add 3 USD per ticket. Variable cost per ticket is 18 USD for walking tours and 34 USD for minibus tours. Total fixed seasonal costs, for either tour type, are 585 USD with plentiful bookings and 675 USD with limited bookings. These fixed costs are paid once per season, and there are no other revenues or costs.

Characterize the key elements of the decision-making situation \sscore{C1-3} 
and construct the payoff table. \sscore{C3-6}
\end{minipage}\par

\begin{minipage}{\linewidth}
\SubsectionBox{Furniture Workshop \score{C4-\sml{2}}}
A furniture workshop must select a production plan before furniture orders are known. A design licence earns a fixed royalty regardless of the order state. The possible states, in column order, are Surging orders, Steady orders, Falling orders. Profits are in thousand USD. No probabilities are available.
\begin{center}\small\setlength{\tabcolsep}{4pt}
\begin{tabular}{@{}lrrr@{}}\toprule
Alternative&$S_1$&$S_2$&$S_3$\\\midrule
Custom benches&104&48&-18\\
Modular desks&146&71&23\\
Repair service&78&56&37\\
Design licence&42&42&42\\\bottomrule
\end{tabular}\end{center}

What would the optimistic decision be? What would the conservative decision be? Justify your answer. \score{2}
\end{minipage}\par

\begin{minipage}{\linewidth}
\SubsectionBox{Water Filter Distribution \score{C4-\sml{2}}}
A water filter supplier must select a distribution programme before regional purchasing is known. A maintenance contract pays a fixed profit under every purchasing state. The possible states, in column order, are Large tenders, Regular tenders, Few tenders. Profits are in thousand USD. No probabilities are available.
\begin{center}\small\setlength{\tabcolsep}{4pt}
\begin{tabular}{@{}lrrr@{}}\toprule
Alternative&$S_1$&$S_2$&$S_3$\\\midrule
Retail network&128&46&-31\\
Clinic supply&96&75&48\\
Tender campaign&182&62&21\\
Maintenance contract&39&39&39\\\bottomrule
\end{tabular}\end{center}

What would the optimistic decision be? What would the conservative decision be? Justify your answer. \score{2}
\end{minipage}\par

\begin{minipage}{\linewidth}
\SubsectionBox{Cold Storage Capacity \score{C4-\sml{2}}}
A food distributor must choose a cold storage plan before warehouse occupancy is known. The possible states, in column order, are Full occupancy, High occupancy, Moderate occupancy, Low occupancy. Profits are in thousand USD. No probabilities are available.
\begin{center}\small\setlength{\tabcolsep}{4pt}
\begin{tabular}{@{}lrrrr@{}}\toprule
Alternative&S1&S2&S3&S4\\\midrule
Own warehouse&142&78&26&-34\\
Regional hub&184&93&46&14\\
Mobile coolers&92&74&58&39\\
Shared storage&67&59&44&28\\\bottomrule
\end{tabular}\end{center}

What would the optimistic decision be? What would the conservative decision be? Justify your answer. \score{2}
\end{minipage}\par

\begin{minipage}{\linewidth}
\SubsectionBox{Museum Gift Stock \score{C5-\sml{2}}}
A museum shop must select a gift stocking plan before visitor spending is known. The possible states, in column order, are High spending, Medium spending, Low spending. Profits are in thousand USD. No probabilities are available. All payoffs are to be maximised.
\begin{center}\small\setlength{\tabcolsep}{4pt}
\begin{tabular}{@{}lrrr@{}}\toprule
Alternative&$S_1$&$S_2$&$S_3$\\\midrule
Collector gifts&92&46&-22\\
Mixed gifts&73&68&18\\
Practical gifts&79&63&41\\\bottomrule
\end{tabular}\end{center}

What decision would be selected using the Minimax Regret criterion? Justify your answer. \score{2}
\end{minipage}\par

\begin{minipage}{\linewidth}
\SubsectionBox{Parcel Sorting Capacity \score{C5-\sml{2}}}
A parcel carrier must select a sorting capacity plan before shipment volumes are known. A handling contract provides a fixed profit under every shipment state. The possible states, in column order, are Shipment surge, Normal shipments, Shipment decline. Profits are in thousand USD. No probabilities are available. All payoffs are to be maximised.
\begin{center}\small\setlength{\tabcolsep}{4pt}
\begin{tabular}{@{}lrrr@{}}\toprule
Alternative&$S_1$&$S_2$&$S_3$\\\midrule
New sorting depot&139&62&-24\\
Expand conveyors&97&82&27\\
Modular sorting&113&74&49\\
Handling contract&45&45&45\\\bottomrule
\end{tabular}\end{center}

What decision would be selected using the Minimax Regret criterion? Justify your answer. \score{2}
\end{minipage}\par

\begin{minipage}{\linewidth}
\SubsectionBox{Bakery Production \score{C5-\sml{2}}}
A bakery must choose its daily bread production plan before neighbourhood purchases are known. The possible states, in column order, are Record purchases, Strong purchases, Typical purchases, Weak purchases. Profits are in USD. No probabilities are available. All payoffs are to be maximised.
\begin{center}\small\setlength{\tabcolsep}{4pt}
\begin{tabular}{@{}lrrrr@{}}\toprule
Alternative&S1&S2&S3&S4\\\midrule
Morning bake&1240&870&590&310\\
Split bake&990&920&650&180\\
Preorder bake&720&680&540&440\\\bottomrule
\end{tabular}\end{center}

What decision would be selected using the Minimax Regret criterion? Justify your answer. \score{2}
\end{minipage}\par

\begin{minipage}{\linewidth}
\SubsectionBox{Ceramic Exports \score{C6-\sml{2}}}
A ceramic producer must choose an export product line before overseas orders are known. The possible states, in column order, are Expanding orders, Stable orders, Contracting orders. Profits are in thousand USD.
\begin{center}\small\setlength{\tabcolsep}{4pt}
\begin{tabular}{@{}lrrr@{}}\toprule
Alternative&$S_1$&$S_2$&$S_3$\\
$P(S_j)$&0.50&0.30&0.20\\\midrule
Artisan series&126&58&-12\\
Tableware series&83&71&32\\
Utility series&60&55&46\\\bottomrule
\end{tabular}\end{center}

What decision would be selected using the EMV or Bayes criterion? Justify your answer. \score{2}
\end{minipage}\par

\begin{minipage}{\linewidth}
\SubsectionBox{Campsite Facilities \score{C6-\sml{2}}}
A campsite must select guest facilities before holiday reservations are known. A concession agreement earns a fixed profit under all reservation states. The possible states, in column order, are Many reservations, Normal reservations, Few reservations. Profits are in thousand USD.
\begin{center}\small\setlength{\tabcolsep}{4pt}
\begin{tabular}{@{}lrrr@{}}\toprule
Alternative&$S_1$&$S_2$&$S_3$\\
$P(S_j)$&0.20&0.35&0.45\\\midrule
Cabin complex&131&49&-21\\
Tent platforms&95&72&28\\
Basic pitches&74&65&52\\
Concession agreement&37&37&37\\\bottomrule
\end{tabular}\end{center}

What decision would be selected using the EMV or Bayes criterion? Justify your answer. \score{2}
\end{minipage}\par

\begin{minipage}{\linewidth}
\SubsectionBox{Audio Production Capacity \score{C6-\sml{2}}}
An audio studio must choose a recording capacity plan before project bookings are known. The possible states, in column order, are Booking surge, Strong bookings, Regular bookings, Scarce bookings. Profits are in thousand USD.
\begin{center}\small\setlength{\tabcolsep}{4pt}
\begin{tabular}{@{}lrrrr@{}}\toprule
Alternative&S1&S2&S3&S4\\
$P(S_j)$&0.10&0.25&0.30&0.35\\\midrule
Full studio team&148&92&34&-37\\
Team plus contractors&112&94&62&21\\
Resident team&79&76&69&55\\\bottomrule
\end{tabular}\end{center}

What decision would be selected using the EMV or Bayes criterion? Justify your answer. \score{2}
\end{minipage}\par


\begin{minipage}{\linewidth}
\SubsectionBox{Irrigation Technology \score{C7-\sml{3}}}
A farming cooperative must choose soil sensors or reservoir pumps before rainfall is known. The possible states, in column order, are Severe drought, Intermittent rain, Sustained rain. Profits are in thousand USD.
\begin{center}\small\setlength{\tabcolsep}{4pt}
\begin{tabular}{@{}lrrr@{}}\toprule
Alternative&$S_1$&$S_2$&$S_3$\\\midrule
Soil sensors ($A_1$)&62&48&54\\
Reservoir pumps ($A_2$)&102&70&18\\\bottomrule
\end{tabular}\end{center}
Suppose $P(S_1)=p$ and $P(S_2)=p/2$; determine $P(S_3)$ from the sum-to-one condition.

The stated parameter range is $0\le p\le1$. Identify the feasible probability domain and all infeasible values in this stated range. Form both EMVs, solve their equality, and state the preferred alternative on each feasible interval and the tie at the threshold. \score{3}
\end{minipage}\par

\begin{minipage}{\linewidth}
\SubsectionBox{Island Freight Contracts \score{C7-\sml{3}}}
An island freight company must choose flexible charters or a scheduled route before port access is known. The possible states, in column order, are Open ports, Restricted access, Frequent closures. Profits are in thousand USD.
\begin{center}\small\setlength{\tabcolsep}{4pt}
\begin{tabular}{@{}lrrr@{}}\toprule
Alternative&$S_1$&$S_2$&$S_3$\\\midrule
Flexible charters ($A_1$)&96&68&22\\
Scheduled route ($A_2$)&44&50&58\\\bottomrule
\end{tabular}\end{center}
Suppose $P(S_1)=2p$ and $P(S_2)=p$; determine $P(S_3)$ from the sum-to-one condition.

The stated parameter range is $0\le p\le1$. Identify the feasible probability domain and all infeasible values in this stated range. Form both EMVs, solve their equality, and state the preferred alternative on each feasible interval and the tie at the threshold. \score{3}
\end{minipage}\par
\end{multicols}
\end{document}
